\documentclass[11pt]{article}
\usepackage[a4paper,margin=1in]{geometry}
\usepackage{fontspec}
\usepackage[boldfont=false]{xeCJK}
\setCJKmainfont{FandolSong-Regular.otf}
\usepackage{amsmath}
\usepackage{amssymb}
\usepackage{array}
\usepackage{booktabs}
\usepackage{graphicx}
\usepackage{adjustbox}
\usepackage{enumitem}
\setlist[itemize]{topsep=2pt,partopsep=0pt,parsep=0pt,itemsep=2pt,leftmargin=1.4em}
\setlist[enumerate]{topsep=2pt,partopsep=0pt,parsep=0pt,itemsep=2pt,leftmargin=1.6em}
\usepackage{parskip}
\usepackage{fvextra}
\fvset{breaklines=true,breakanywhere=true,fontsize=\small}
\setlength{\parindent}{0pt}

\begin{document}

\begin{center}
  {\LARGE\bfseries Topic 07}\\[0.35em]
  {\large Igcse Computer Science}
\end{center}
\vspace{0.6em}

\section*{The program development life cycle}

The \textbf{program development life cycle} 程序开发生命周期 is the set of stages used to make a program. There are four stages.

\begin{center}
\adjustbox{max width=\linewidth}{%
\begin{tabular}{ll}
\toprule
\textbf{Stage} & \textbf{What you do} \\
\midrule
\textbf{analysis} 分析 & study the problem and work out what is needed \\
\textbf{design} 设计 & plan how the program will work \\
\textbf{coding} 编码 & write the program code and test it as you go \\
\textbf{testing} 测试 & run the finished program with test data to find errors \\
\bottomrule
\end{tabular}}
\end{center}

\subsection*{Analysis}

In analysis you understand the problem. Two key skills help:

\begin{itemize}
\item \textbf{abstraction} 抽象 — keep only the important details and ignore the rest;

\item \textbf{decomposition} 分解 — break a big problem into smaller, easier parts.

\end{itemize}

\subsection*{Design}

In design you plan the solution, often using decomposition. You can show the parts as \textbf{sub-systems} 子系统 in a \textbf{structure diagram} 结构图 (a chart that splits a system into smaller boxes).

\subsection*{Coding and testing}

In \textbf{coding} you write the program code. You use \textbf{iterative testing} 迭代测试 — test small parts again and again as you build them. In \textbf{testing} you run the whole program with \textbf{test data} 测试数据 to check it works.

\section*{Design tools}

You can plan a solution in three main ways.

\begin{itemize}
\item a \textbf{structure diagram} — shows the parts of a system and how they fit together;

\item a \textbf{flowchart} 流程图 — a diagram using boxes and arrows to show the steps in order;

\item \textbf{pseudocode} 伪代码 — steps written in simple, code-like English (not a real language).

\end{itemize}

\section*{Algorithms}

An \textbf{algorithm} 算法 is a set of steps, in the right order, that solves a problem. Every algorithm can be split into three parts:

\begin{itemize}
\item \textbf{input} 输入 — the data that goes in;

\item \textbf{processing} 处理 — the work done on the data;

\item \textbf{output} 输出 — the result that comes out.

\end{itemize}

This is called decomposition into inputs, processes and outputs. For example, for "find the average of three marks": the inputs are the three marks; the processing is adding them and dividing by 3; the output is the average.

\section*{Validation and verification}

When data is entered, you check it to reduce mistakes.

\textbf{Validation} 验证 checks that the data is \textbf{sensible} and follows the rules. It cannot check that the data is true, only that it is allowed.

\begin{center}
\adjustbox{max width=\linewidth}{%
\begin{tabular}{ll}
\toprule
\textbf{Validation check} & \textbf{What it checks} \\
\midrule
\textbf{range check} 范围检查 & the value is between a lowest and highest allowed value \\
\textbf{length check} 长度检查 & the number of characters is allowed (e.g. a password ≥ 8) \\
\textbf{type check} 类型检查 & the data is the right type (e.g. a number, not letters) \\
\textbf{presence check} 存在性检查 & something has actually been entered (not left blank) \\
\textbf{format check} 格式检查 & the data is in the right pattern (e.g. a date as dd/mm/yyyy) \\
\textbf{check digit} 校验码 & an extra digit confirms a number was entered correctly \\
\bottomrule
\end{tabular}}
\end{center}

\textbf{Verification} 核实 checks that data was \textbf{copied or entered correctly} (no mistakes while typing it in). Two methods:

\begin{itemize}
\item \textbf{visual check} 目视检查 — a person compares the typed data with the original;

\item \textbf{double entry} 双重输入 — the data is entered twice and the two copies are compared.

\end{itemize}

\section*{Trace tables}

A \textbf{trace table} 追踪表 records the value of each variable as an algorithm runs, step by step. It helps you:

\begin{itemize}
\item check that an algorithm works correctly;

\item work out \textbf{what an algorithm does} by following it with given data.

\end{itemize}

Example: trace this algorithm with the input 5.

\begin{Verbatim}
INPUT n
total ← 0
FOR i ← 1 TO n
    total ← total + i
NEXT i
OUTPUT total
\end{Verbatim}

\begin{center}
\adjustbox{max width=\linewidth}{%
\begin{tabular}{lll}
\toprule
\textbf{i} & \textbf{total} & \textbf{OUTPUT} \\
\midrule
1 & 1 &  \\
2 & 3 &  \\
3 & 6 &  \\
4 & 10 &  \\
5 & 15 & 15 \\
\bottomrule
\end{tabular}}
\end{center}

The trace shows the algorithm adds up 1 to n. With input 5 the output is 15.

\section*{Test data}

\textbf{Test data} is data you use to test a program. There are four types you must know.

\begin{center}
\adjustbox{max width=\linewidth}{%
\begin{tabular}{lll}
\toprule
\textbf{Type} & \textbf{Meaning} & \textbf{Example (age 0–120 allowed)} \\
\midrule
\textbf{normal} 正常数据 & sensible data that should be accepted & 25 \\
\textbf{abnormal} 异常数据 & wrong data that should be rejected & -4 or "cat" \\
\textbf{extreme} 极端数据 & the largest and smallest values still allowed & 0 and 120 \\
\textbf{boundary} 边界数据 & the values on each side of a limit (one allowed, one not) & 120 and 121 \\
\bottomrule
\end{tabular}}
\end{center}

\section*{Standard methods of solution}

You must know these common algorithms.

\subsection*{Linear search}

A \textbf{linear search} 线性查找 checks each item in a list, one by one, until it finds the value it wants or reaches the end.

\begin{Verbatim}
found ← FALSE
FOR i ← 0 TO 9
    IF list[i] = searchValue THEN
        found ← TRUE
    ENDIF
NEXT i
OUTPUT found
\end{Verbatim}

\subsection*{Bubble sort}

A \textbf{bubble sort} 冒泡排序 puts a list in order. It compares each pair of side-by-side items and swaps them if they are in the wrong order. It repeats this until no more swaps are needed.

\begin{Verbatim}
FOR i ← 0 TO 8
    IF list[i] > list[i + 1] THEN
        temp ← list[i]
        list[i] ← list[i + 1]
        list[i + 1] ← temp
    ENDIF
NEXT i
\end{Verbatim}

\subsection*{Totalling and counting}

\begin{itemize}
\item \textbf{totalling} 求和 — keep adding values to a running total (\texttt{total ← total + value}).

\item \textbf{counting} 计数 — add 1 to a counter each time something happens (\texttt{count ← count + 1}).

\end{itemize}

\subsection*{Maximum, minimum and average}

\begin{itemize}
\item to find the \textbf{maximum} 最大值: keep the largest value seen so far.

\item to find the \textbf{minimum} 最小值: keep the smallest value seen so far.

\item to find the \textbf{average} 平均值: divide the total by how many values there are.

\end{itemize}

\begin{Verbatim}
total ← 0
FOR i ← 0 TO 9
    total ← total + list[i]
NEXT i
average ← total / 10
OUTPUT average
\end{Verbatim}


\end{document}
